El gas radó és una de les fonts de radioactivitat natural més abundants de la Terra. El radó prové de la descomposició d'elements radioactius naturals (com l'urani i el tori). El gas es difon a través del sòl fins a arribar a la superfície. La cadena de desintegració del radó $^{222}_{86}\mathrm{Rn}$ inclou vuit desintegracions radioactives, fins que es forma l'isòtop estable del plom $^{206}_{82}\mathrm{Pb}$. En la figura següent es representen els nuclis que formen part d'aquesta cadena de desintegració nuclear. Al costat de cada nucli, se n'indica el període de semidesintegració.

\figura[0.75]{fig1.png}

\begin{apartat}{1,25}
Escriviu les reaccions nuclears que permeten arribar al $^{206}_{82}\mathrm{Pb}$ a partir del $^{222}_{86}\mathrm{Rn}$.
\end{apartat}

\begin{apartat}{1,25}
El gràfic següent correspon a l'evolució dels nuclis d'una de les desintegracions radioactives de la cadena del radó. La mostra estudiada inicialment tenia $8,00\times 10^{23}$ nuclis. A partir del gràfic, determineu quin és el període de semidesintegració de la mostra, i raoneu a quin nucli de la cadena correspon. Amb aquesta dada calculeu quants dies han de passar fins que s'hagin desintegrat $7,95\times 10^{23}$ àtoms.

\figura[0.5]{fig2.png}
\end{apartat}
